\documentclass[12pt]{ctexart}
\usepackage{amsmath,amssymb,bm}
\usepackage[a4paper,margin=1in]{geometry}
\usepackage{physics}
\usepackage{mathtools}

\begin{document}

\section*{Kirchhoff--Love 薄板振动讲义（扩充版）}

\subsection*{1. 统一框架：线性薄板方程、模态与静/动态分解}

\subsubsection*{1.1 基本假设与控制方程}
考虑厚度为 \(h\) 的均匀、各向同性、线弹性薄板，其中面区域记为 \(\Omega\)，边界记为 \(\Gamma\)。采用 Kirchhoff--Love 假设：
\begin{itemize}
\item 板厚远小于面内尺度；
\item 挠度与转角都足够小；
\item 变形后中面法线仍保持为直线，并保持与中面垂直；
\item 忽略横向剪切变形与转动惯量。
\end{itemize}

记中面横向位移为 \(w(x,y,t)\)，则板的控制方程为
\[
\rho h \frac{\partial^2 w}{\partial t^2} + D \nabla_t^4 w = q(x,y,t),
\qquad
D=\frac{Eh^3}{12(1-\nu^2)},
\]
其中
\[
\nabla_t^2=\frac{\partial^2}{\partial x^2}+\frac{\partial^2}{\partial y^2},
\qquad
\nabla_t^4=(\nabla_t^2)^2.
\]

为了统一静态、自由振动与受迫振动，先考虑谐振型载荷
\[
q(x,y,t)=\Re\{Q(x,y)e^{-i\Omega t}\},
\qquad
w(x,y,t)=\Re\{W(x,y)e^{-i\Omega t}\}.
\]
代入得
\[
\left(D\nabla_t^4-\rho h \Omega^2\right)W=Q.
\]
于是三类问题都可看成这一统一方程的特例：
\[
\begin{cases}
D\nabla_t^4 w_s=q_s, & \text{静态问题，对应 } \Omega=0,\\[4pt]
\left(D\nabla_t^4-\rho h \omega^2\right)\Phi=0, & \text{自由振动，对应 } Q=0,\\[4pt]
\left(D\nabla_t^4-\rho h \Omega^2\right)W=Q, & \text{谐受迫振动。}
\end{cases}
\]

\subsubsection*{1.2 弯矩、剪力与边界条件}
记板的弯矩张量
\[
\bm M=-D\left(\nu \bm I_t \nabla_t^2 +(1-\nu)\bm\nabla_t\bm\nabla_t\right)w.
\]
在直角坐标下，
\[
M_x=-D\left(\frac{\partial^2 w}{\partial x^2}+\nu\frac{\partial^2 w}{\partial y^2}\right),
\qquad
M_y=-D\left(\frac{\partial^2 w}{\partial y^2}+\nu\frac{\partial^2 w}{\partial x^2}\right),
\]
\[
M_{xy}=-D(1-\nu)\frac{\partial^2 w}{\partial x\partial y}.
\]
相应的 Kirchhoff 剪力为
\[
Q_x=-D\frac{\partial}{\partial x}\left(\nabla_t^2 w\right),
\qquad
Q_y=-D\frac{\partial}{\partial y}\left(\nabla_t^2 w\right).
\]
边界上的有效法向剪力记为
\[
V_n = Q_n+\frac{\partial M_{n\tau}}{\partial \tau}.
\]

常见三类边界条件统一写为
\[
\begin{cases}
w|_\Gamma=0,\qquad \dfrac{\partial w}{\partial n}\bigg|_\Gamma=0, & \text{C (固支)},\\[8pt]
w|_\Gamma=0,\qquad M_{nn}|_\Gamma=0, & \text{S (简支)},\\[8pt]
M_{nn}|_\Gamma=0,\qquad V_n|_\Gamma=0, & \text{F (自由)}.
\end{cases}
\]

\subsubsection*{1.3 本征模态与正交性}
自由振动满足
\[
D\nabla_t^4\Phi_n=\rho h \omega_n^2 \Phi_n,
\]
并在边界上满足与原问题相同的齐次边界条件。对经典的 C / S / F 边界条件，该本征值问题是自伴的，因此不同本征频率对应的模态满足正交关系
\[
\int_\Omega \rho h\, \Phi_m \Phi_n \, d\Omega = 0,
\qquad m\neq n.
\]
记模态质量为
\[
M_n=\int_\Omega \rho h\, \Phi_n^2\, d\Omega.
\]
则可将模态选为质量正交归一基
\[
\int_\Omega \rho h\, \Phi_m \Phi_n\, d\Omega = M_n\delta_{mn}.
\]

\paragraph{注 1}
若某个本征频率具有重根，则可在对应本征子空间中取一组相互正交的基函数。圆板中 \(m\ge 1\) 的 \(\cos m\phi\) 与 \(\sin m\phi\) 正是这种二重简并的典型例子。

\paragraph{注 2}
若整块板四周全自由，则存在刚体模态，对应 \(\omega=0\)。这时静力问题若无额外约束，解一般不唯一。

\subsubsection*{1.4 模态展开：静态、自由振动与受迫振动的统一}
设
\[
w(x,y,t)=\sum_{n=1}^\infty q_n(t)\Phi_n(x,y).
\]
代入控制方程并利用正交性，得到一组完全解耦的单自由度方程
\[
\ddot q_n(t)+\omega_n^2 q_n(t)=\frac{F_n(t)}{M_n},
\]
其中广义力为
\[
F_n(t)=\int_\Omega q(x,y,t)\Phi_n(x,y)\, d\Omega.
\]

于是：

\paragraph{(i) 静态问题}
若载荷与时间无关，且静态解写作
\[
w_s(x,y)=\sum_{n=1}^\infty a_n\Phi_n(x,y),
\]
则
\[
a_n=\frac{1}{M_n\omega_n^2}\int_\Omega q_s(x,y)\Phi_n(x,y)\, d\Omega,
\]
故
\[
w_s(x,y)=
\sum_{n=1}^\infty
\frac{\displaystyle \int_\Omega q_s \Phi_n\, d\Omega}
{M_n\omega_n^2}
\Phi_n(x,y).
\]

\paragraph{(ii) 自由振动}
对应 \(F_n=0\)，故
\[
q_n(t)=A_n\cos\omega_n t+B_n\sin\omega_n t.
\]
于是
\[
w(x,y,t)=\sum_{n=1}^\infty \left(A_n\cos\omega_n t+B_n\sin\omega_n t\right)\Phi_n(x,y).
\]

\paragraph{(iii) 谐受迫振动}
若
\[
q(x,y,t)=\Re\{Q(x,y)e^{-i\Omega t}\},
\qquad
w(x,y,t)=\Re\{W(x,y)e^{-i\Omega t}\},
\]
且写
\[
W(x,y)=\sum_{n=1}^\infty c_n \Phi_n(x,y),
\]
则
\[
c_n=
\frac{\displaystyle \int_\Omega Q\Phi_n\, d\Omega}
{M_n(\omega_n^2-\Omega^2)}.
\]
故谐响应为
\[
W(x,y)=
\sum_{n=1}^\infty
\frac{\displaystyle \int_\Omega Q\Phi_n\, d\Omega}
{M_n(\omega_n^2-\Omega^2)}
\Phi_n(x,y).
\]
由此可见，静态问题就是 \(\Omega=0\) 的极限；自由振动则是 \(Q=0\) 且只保留齐次解。

\subsection*{2. 圆板}

考虑半径为 \(R\) 的实心圆板，采用极坐标 \((r,\phi)\)，并记
\[
\xi=\frac{r}{R}.
\]

\subsubsection*{2.1 极坐标下的边界算子}
极坐标下
\[
\nabla_t^2
=
\frac{\partial^2}{\partial r^2}
+\frac{1}{r}\frac{\partial}{\partial r}
+\frac{1}{r^2}\frac{\partial^2}{\partial \phi^2}.
\]
弯矩分量为
\[
M_{rr}
=
-D\left(
\frac{\partial^2 w}{\partial r^2}
+\nu\left(
\frac{1}{r}\frac{\partial w}{\partial r}
+\frac{1}{r^2}\frac{\partial^2 w}{\partial \phi^2}
\right)\right),
\]
\[
M_{\phi\phi}
=
-D\left(
\frac{1}{r}\frac{\partial w}{\partial r}
+\frac{1}{r^2}\frac{\partial^2 w}{\partial \phi^2}
+\nu \frac{\partial^2 w}{\partial r^2}
\right),
\]
\[
M_{r\phi}
=
-D(1-\nu)\left(
\frac{1}{r}\frac{\partial^2 w}{\partial r\partial \phi}
-\frac{1}{r^2}\frac{\partial w}{\partial \phi}
\right).
\]
有效径向剪力为
\[
V_r
=
-D\left[
\frac{\partial}{\partial r}\left(\nabla_t^2 w\right)
+
(1-\nu)\left(
\frac{1}{r^2}\frac{\partial^3 w}{\partial r\partial \phi^2}
-\frac{1}{r^3}\frac{\partial^2 w}{\partial \phi^2}
\right)
\right].
\]
因此 \(r=R\) 上的三类边界条件为
\[
\begin{cases}
w=0,\qquad \dfrac{\partial w}{\partial r}=0, & \text{C},\\[8pt]
w=0,\qquad M_{rr}=0, & \text{S},\\[8pt]
M_{rr}=0,\qquad V_r=0, & \text{F}.
\end{cases}
\]

\subsubsection*{2.2 统一的分离变量形式}
对于谐响应
\[
\left(\nabla_t^4-k^4\right)W=\frac{Q(r,\phi)}{D},
\qquad
k^4=\frac{\rho h \Omega^2}{D},
\]
将角向展开为
\[
W(r,\phi)=\sum_{m=0}^\infty W_m(r)\cos(m\phi)
+\sum_{m=1}^\infty \widetilde W_m(r)\sin(m\phi).
\]
对于每个角向阶数 \(m\)，齐次方程的径向通解为
\[
W_m^{(h)}(r)
=
A_m J_m(kr)+B_m Y_m(kr)+C_m I_m(kr)+D_m K_m(kr).
\]
对实心圆板，由于 \(r=0\) 处需有限，必须舍去 \(Y_m,K_m\)，故
\[
W_m^{(h)}(r)=A_m J_m(kr)+B_m I_m(kr).
\]

静态问题是 \(k=0\) 的退化极限。此时双调和方程的齐次解为：

\paragraph{(i) \(m=0\)}
\[
w_0^{(h)}(\xi)=A_0+B_0\xi^2+C_0\ln\xi+D_0\xi^2\ln\xi.
\]

\paragraph{(ii) \(m=1\)}
\[
w_1^{(h)}(\xi)=A_1\xi+B_1\xi^{-1}+C_1\xi^3+D_1\xi\ln\xi.
\]

\paragraph{(iii) \(m\ge 2\)}
\[
w_m^{(h)}(\xi)=A_m\xi^m+B_m\xi^{-m}+C_m\xi^{m+2}+D_m\xi^{2-m}.
\]

对于实心圆板，由 \(\xi=0\) 处有限性知：
\[
w_0^{(h)}(\xi)=A_0+B_0\xi^2,
\]
\[
w_1^{(h)}(\xi)=A_1\xi+C_1\xi^3,
\]
\[
w_m^{(h)}(\xi)=A_m\xi^m+C_m\xi^{m+2},\qquad m\ge 2.
\]

\subsubsection*{2.3 圆板的模态}
自由振动时 \(Q=0,\ \Omega=\omega\)，记
\[
\beta=kR=\omega^{1/2}R\left(\frac{\rho h}{D}\right)^{1/4},
\qquad
\Lambda=\beta^2=\omega R^2\sqrt{\frac{\rho h}{D}}.
\]
于是实心圆板的自由振型可写为
\[
\Phi_{m,s}^{(c)}(r,\phi)
=
\left[A J_m(\beta_{m,s}\xi)+B I_m(\beta_{m,s}\xi)\right]\cos(m\phi),
\]
\[
\Phi_{m,s}^{(s)}(r,\phi)
=
\left[A J_m(\beta_{m,s}\xi)+B I_m(\beta_{m,s}\xi)\right]\sin(m\phi),
\]
其中 \(s=1,2,\cdots\) 表示同一 \(m\) 下按径向节点数排序的模态序号。对应固有频率为
\[
\omega_{m,s}
=
\frac{\beta_{m,s}^2}{R^2}\sqrt{\frac{D}{\rho h}}.
\]
当 \(m=0\) 时模态非简并；当 \(m\ge 1\) 时 \(\cos(m\phi)\) 与 \(\sin(m\phi)\) 为同频二重模态。

\subsubsection*{2.4 实心圆板三类边界条件下的本征方程}

\paragraph{(i) 固支边界 C}
边界条件
\[
W(1,\phi)=0,\qquad \pdv{W}{\xi}(1,\phi)=0.
\]
对应系数方程
\[
\begin{pmatrix}
J_m(\beta) & I_m(\beta)\\
J_m'(\beta) & I_m'(\beta)
\end{pmatrix}
\begin{pmatrix} A\\ B \end{pmatrix}=0.
\]
非平凡解存在的条件为
\[
J_m(\beta)I_m'(\beta)-J_m'(\beta)I_m(\beta)=0,
\]
进一步用递推关系可化为
\[
J_m(\beta) I_{m+1}(\beta)+I_m(\beta) J_{m+1}(\beta)=0.
\]
因此固支圆板的最终结果为：
\[
\beta=\beta_{m,s}\ \text{由}\ 
J_m(\beta) I_{m+1}(\beta)+I_m(\beta) J_{m+1}(\beta)=0
\ \text{给出},
\]
\[
\omega_{m,s}
=
\frac{\beta_{m,s}^2}{R^2}\sqrt{\frac{D}{\rho h}},
\]
\[
\Phi_{m,s}(r,\phi)
=
\left[
J_m(\beta_{m,s}\xi)
-\frac{J_m(\beta_{m,s})}{I_m(\beta_{m,s})}I_m(\beta_{m,s}\xi)
\right]
\begin{cases}
\cos(m\phi),\\
\sin(m\phi).
\end{cases}
\]

\paragraph{(ii) 简支边界 S}
边界条件
\[
W(1,\phi)=0,\qquad M_{rr}(1,\phi)=0.
\]
对应线性方程组
\[
\begin{pmatrix}
J_m(\beta) & I_m(\beta)\\
\beta^2 J_m(\beta)+(1-\nu)\bigl(\beta J_m'(\beta)-m^2J_m(\beta)\bigr)
&
-\beta^2 I_m(\beta)+(1-\nu)\bigl(\beta I_m'(\beta)-m^2I_m(\beta)\bigr)
\end{pmatrix}
\begin{pmatrix} A\\ B \end{pmatrix}=0.
\]
本征方程可化简为
\[
\frac{J_{m+1}(\beta)}{J_m(\beta)}+
\frac{I_{m+1}(\beta)}{I_m(\beta)}
=
\frac{2\beta}{1-\nu}.
\]
因此简支圆板的最终结果为：
\[
\beta=\beta_{m,s}\ \text{由}\ 
\frac{J_{m+1}(\beta)}{J_m(\beta)}+
\frac{I_{m+1}(\beta)}{I_m(\beta)}
=
\frac{2\beta}{1-\nu}
\ \text{给出},
\]
\[
\omega_{m,s}
=
\frac{\beta_{m,s}^2}{R^2}\sqrt{\frac{D}{\rho h}},
\]
\[
\Phi_{m,s}(r,\phi)
=
\left[
J_m(\beta_{m,s}\xi)
-\frac{J_m(\beta_{m,s})}{I_m(\beta_{m,s})}I_m(\beta_{m,s}\xi)
\right]
\begin{cases}
\cos(m\phi),\\
\sin(m\phi).
\end{cases}
\]

\paragraph{(iii) 自由边界 F}
边界条件
\[
M_{rr}(1,\phi)=0,\qquad V_r(1,\phi)=0.
\]
对应方程组
\[
\begin{pmatrix}
\beta^2 J_m(\beta)+(1-\nu)\bigl(\beta J_m'(\beta)-m^2J_m(\beta)\bigr)
&
-\beta^2 I_m(\beta)+(1-\nu)\bigl(\beta I_m'(\beta)-m^2I_m(\beta)\bigr)
\\[6pt]
\beta^3 J_m'(\beta)+m^2(1-\nu)\bigl(\beta J_m'(\beta)-J_m(\beta)\bigr)
&
-\beta^3 I_m'(\beta)+m^2(1-\nu)\bigl(\beta I_m'(\beta)-I_m(\beta)\bigr)
\end{pmatrix}
\begin{pmatrix} A\\ B \end{pmatrix}=0.
\]
因此自由圆板的最终结果为：
\[
\beta=\beta_{m,s}\ \text{由上式 }2\times2\text{ 行列式为零确定},
\]
\[
\omega_{m,s}
=
\frac{\beta_{m,s}^2}{R^2}\sqrt{\frac{D}{\rho h}}.
\]
若取
\[
a_m(\beta)
=
\beta^2 J_m(\beta)+(1-\nu)\bigl(\beta J_m'(\beta)-m^2J_m(\beta)\bigr),
\]
\[
b_m(\beta)
=
-\beta^2 I_m(\beta)+(1-\nu)\bigl(\beta I_m'(\beta)-m^2I_m(\beta)\bigr),
\]
则对应振型可写为
\[
\Phi_{m,s}(r,\phi)
=
\left[
J_m(\beta_{m,s}\xi)-\frac{a_m(\beta_{m,s})}{b_m(\beta_{m,s})}I_m(\beta_{m,s}\xi)
\right]
\begin{cases}
\cos(m\phi),\\
\sin(m\phi).
\end{cases}
\]

\subsubsection*{2.5 轴对称静态问题的最终结果}
若载荷轴对称，则只保留 \(m=0\) 项。静态方程化为
\[
\frac{1}{r}\frac{d}{dr}\left[
r\frac{d}{dr}
\left(
\frac{1}{r}\frac{d}{dr}\left(r\frac{dw}{dr}\right)
\right)
\right]
=
\frac{q(r)}{D}.
\]

\paragraph{(i) 均布载荷 \(q=q_0\)}
轴对称特解可取
\[
w^{(p)}=\frac{q_0R^4}{64D}\xi^4.
\]
实心圆板轴对称总解为
\[
w(\xi)=\frac{q_0R^4}{64D}\left(\xi^4+A\xi^2+B\right).
\]

\textbf{固支边界 C：}
\[
w(1)=0,\qquad \left.\dv{w}{\xi}\right|_{\xi=1}=0
\]
给出
\[
w(\xi)=\frac{q_0R^4}{64D}(1-\xi^2)^2.
\]

\textbf{简支边界 S：}
\[
w(1)=0,\qquad M_{rr}(1)=0
\]
给出
\[
w(\xi)=\frac{q_0R^4}{64D}
\left(
\xi^4-\frac{2(3+\nu)}{1+\nu}\xi^2+\frac{5+\nu}{1+\nu}
\right).
\]

\textbf{自由边界 F：}
整体自由的静力问题存在刚体模态，若无额外支承与约束，则静平衡解一般不唯一，因此此处不作单独讨论。

\paragraph{(ii) 中心点力 \(P\)}
对作用于圆心的集中力，轴对称解在 \(r>0\) 内满足齐次双调和方程，其最终结果为：

\textbf{固支边界 C：}
\[
w(\xi)
=
\frac{PR^2}{16\pi D}
\left(
1-\xi^2+2\xi^2\ln\xi
\right).
\]

\textbf{简支边界 S：}
\[
w(\xi)
=
\frac{PR^2}{16\pi D}
\left[
\frac{3+\nu}{1+\nu}(1-\xi^2)+2\xi^2\ln\xi
\right].
\]

这两个公式在 \(\xi\to 0\) 时有限，其导数在圆心附近呈现集中载荷所对应的特征奇性。

\subsubsection*{2.6 非轴对称静载的一般结果}
若静载 \(q(r,\phi)\) 非轴对称，可展开为
\[
q(r,\phi)=q_0(r)+\sum_{m=1}^\infty q_m^{(c)}(r)\cos m\phi
+\sum_{m=1}^\infty q_m^{(s)}(r)\sin m\phi.
\]
对应挠度亦可展开为
\[
w(r,\phi)=w_0(r)+\sum_{m=1}^\infty w_m^{(c)}(r)\cos m\phi
+\sum_{m=1}^\infty w_m^{(s)}(r)\sin m\phi.
\]
每一阶径向函数由前述双调和齐次解加上相应特解构成，系数由
\[
r=0 \text{ 处有限性}+\ r=R \text{ 处边界条件}
\]
唯一确定。故对圆板而言，\textbf{最终结果的通式}就是：
\[
w_m(r)=w_m^{(p)}(r)+w_m^{(h)}(r),
\]
其中 \(w_m^{(h)}\) 的一般形式已经在 \S 2.2 中列出，而 \(w_m^{(p)}\) 则由具体载荷项决定。

\subsection*{3. 长方板}

取长方板区域为
\[
0<x<a,\qquad 0<y<b.
\]

\subsubsection*{3.1 基本方程与边界算子}
谐响应方程为
\[
\left(\nabla_t^4-k^4\right)W=\frac{Q(x,y)}{D},
\qquad
k^4=\frac{\rho h \Omega^2}{D}.
\]
静态问题对应 \(k=0\)：
\[
\nabla_t^4 w=\frac{q(x,y)}{D}.
\]
自由振动对应
\[
\left(\nabla_t^4-k^4\right)\Phi=0.
\]

沿 \(x=\text{const}\) 的边，
\[
M_x=-D\left(W_{xx}+\nu W_{yy}\right),
\qquad
V_x=-D\left(W_{xxx}+(2-\nu)W_{xyy}\right).
\]
沿 \(y=\text{const}\) 的边，
\[
M_y=-D\left(W_{yy}+\nu W_{xx}\right),
\qquad
V_y=-D\left(W_{yyy}+(2-\nu)W_{xxy}\right).
\]
因此：
\[
\begin{cases}
W=0,\qquad \partial_n W=0, & \text{C},\\[6pt]
W=0,\qquad M_{nn}=0, & \text{S},\\[6pt]
M_{nn}=0,\qquad V_n=0, & \text{F}.
\end{cases}
\]

\subsubsection*{3.2 四边简支长方板：Navier 完整结果}
这是最标准、也是最完整可显式写出的情形。设四边均为简支：
\[
W(0,y)=W(a,y)=0,\qquad W(x,0)=W(x,b)=0,
\]
\[
M_x|_{x=0,a}=0,\qquad M_y|_{y=0,b}=0.
\]

\paragraph{(i) 自由振动模态与频率}
取
\[
\Phi_{mn}(x,y)=\sin\frac{m\pi x}{a}\sin\frac{n\pi y}{b},
\qquad m,n=1,2,\cdots
\]
则
\[
\nabla_t^2 \Phi_{mn}
=
-\lambda_{mn}^2 \Phi_{mn},
\qquad
\lambda_{mn}^2=
\left(\frac{m\pi}{a}\right)^2+\left(\frac{n\pi}{b}\right)^2.
\]
进一步
\[
\nabla_t^4 \Phi_{mn}=\lambda_{mn}^4 \Phi_{mn}.
\]
故本征频率为
\[
\omega_{mn}
=
\sqrt{\frac{D}{\rho h}}\lambda_{mn}^2
=
\pi^2\sqrt{\frac{D}{\rho h}}
\left(
\frac{m^2}{a^2}+\frac{n^2}{b^2}
\right).
\]
于是最终结果为
\[
\Phi_{mn}(x,y)=\sin\frac{m\pi x}{a}\sin\frac{n\pi y}{b},
\]
\[
\omega_{mn}
=
\pi^2\sqrt{\frac{D}{\rho h}}
\left(
\frac{m^2}{a^2}+\frac{n^2}{b^2}
\right).
\]

\paragraph{(ii) 任意谐载荷下的显式响应}
将载荷展开为
\[
Q(x,y)=\sum_{m=1}^\infty\sum_{n=1}^\infty Q_{mn}
\sin\frac{m\pi x}{a}\sin\frac{n\pi y}{b},
\]
其中
\[
Q_{mn}
=
\frac{4}{ab}\int_0^a\int_0^b
Q(x,y)\sin\frac{m\pi x}{a}\sin\frac{n\pi y}{b}\, dx\,dy.
\]
则谐响应显式为
\[
W(x,y)=
\sum_{m=1}^\infty\sum_{n=1}^\infty
\frac{Q_{mn}}{D\lambda_{mn}^4-\rho h \Omega^2}
\sin\frac{m\pi x}{a}\sin\frac{n\pi y}{b}.
\]
这就是四边简支板在任意谐载荷下的最终结果。

\paragraph{(iii) 静态问题}
取 \(\Omega=0\)，即 \(Q=q\)，则
\[
w(x,y)=
\sum_{m=1}^\infty\sum_{n=1}^\infty
\frac{Q_{mn}}{D\lambda_{mn}^4}
\sin\frac{m\pi x}{a}\sin\frac{n\pi y}{b}.
\]

\paragraph{(iv) 均布载荷 \(q=q_0\)}
此时
\[
Q_{mn}
=
\frac{16q_0}{mn\pi^2},
\qquad m,n\ \text{均为奇数},
\]
偶数项为零，因此
\[
w(x,y)=
\frac{16q_0}{D\pi^6}
\sum_{\substack{m,n=1\\ m,n\ \text{odd}}}^{\infty}
\frac{
\sin\frac{m\pi x}{a}\sin\frac{n\pi y}{b}
}{
mn\left(\dfrac{m^2}{a^2}+\dfrac{n^2}{b^2}\right)^2
}.
\]

\subsubsection*{3.3 一对对边简支长方板：L\'evy 完整框架}
若 \(x=0,a\) 两边为简支，则
\[
W(0,y)=W(a,y)=0,\qquad M_x|_{x=0,a}=0.
\]
于是可取 L\'evy 型展开
\[
W(x,y)=\sum_{m=1}^\infty Y_m(y)\sin\frac{m\pi x}{a}.
\]
记
\[
\alpha_m=\frac{m\pi}{a}.
\]
则每个 \(Y_m(y)\) 满足
\[
\left(\frac{d^2}{dy^2}-\alpha_m^2\right)^2Y_m-k^4Y_m=\frac{Q_m(y)}{D},
\]
其中
\[
Q_m(y)=\frac{2}{a}\int_0^a Q(x,y)\sin\frac{m\pi x}{a}\, dx.
\]

\paragraph{(i) 自由振动与谐响应的齐次解}
当 \(Q_m=0\) 时，令
\[
p_m^2=\alpha_m^2+k^2,\qquad
q_m^2=k^2-\alpha_m^2,
\]
则齐次解可写为
\[
Y_m^{(h)}(y)
=
A_m\cosh(p_my)+B_m\sinh(p_my)+C_m\cos(q_my)+D_m\sin(q_my).
\]
于是 L\'evy 类长方板的最终结果总是
\[
W(x,y)=
\sum_{m=1}^\infty
\left[
A_m\cosh(p_my)+B_m\sinh(p_my)+C_m\cos(q_my)+D_m\sin(q_my)
\right]
\sin\frac{m\pi x}{a},
\]
其中 \(A_m,B_m,C_m,D_m\) 由 \(y=0,b\) 两边的边界条件决定。

\paragraph{(ii) 便于写边界条件的四个行向量}
定义
\[
R_W(y)=
\begin{bmatrix}
\cosh(p_my) & \sinh(p_my) & \cos(q_my) & \sin(q_my)
\end{bmatrix},
\]
\[
R_\theta(y)=
\begin{bmatrix}
p_m\sinh(p_my) & p_m\cosh(p_my) & -q_m\sin(q_my) & q_m\cos(q_my)
\end{bmatrix},
\]
\[
R_M(y)=
\begin{bmatrix}
(p_m^2-\nu\alpha_m^2)\cosh(p_my) &
(p_m^2-\nu\alpha_m^2)\sinh(p_my) &
-(q_m^2+\nu\alpha_m^2)\cos(q_my) &
-(q_m^2+\nu\alpha_m^2)\sin(q_my)
\end{bmatrix},
\]
\[
R_V(y)=
\begin{bmatrix}
p_m\bigl(p_m^2-(2-\nu)\alpha_m^2\bigr)\sinh(p_my) &
p_m\bigl(p_m^2-(2-\nu)\alpha_m^2\bigr)\cosh(p_my) &
q_m\bigl(q_m^2+(2-\nu)\alpha_m^2\bigr)\sin(q_my) &
-q_m\bigl(q_m^2+(2-\nu)\alpha_m^2\bigr)\cos(q_my)
\end{bmatrix}.
\]

记
\[
\bm c_m=
\begin{bmatrix}
A_m\\ B_m\\ C_m\\ D_m
\end{bmatrix}.
\]
则
\[
Y_m(y)=R_W(y)\bm c_m,\qquad
Y_m'(y)=R_\theta(y)\bm c_m,
\]
\[
Y_m''(y)-\nu\alpha_m^2Y_m(y)=R_M(y)\bm c_m,
\]
\[
Y_m'''(y)-(2-\nu)\alpha_m^2Y_m'(y)=R_V(y)\bm c_m.
\]

\paragraph{(iii) 各类 \(y\)-边界条件的最终特征方程}
在 \(y=0,b\) 上，三类边界条件分别对应：

\[
\text{C at } y=y_0:\qquad
R_W(y_0)\bm c_m=0,\qquad
R_\theta(y_0)\bm c_m=0;
\]
\[
\text{S at } y=y_0:\qquad
R_W(y_0)\bm c_m=0,\qquad
R_M(y_0)\bm c_m=0;
\]
\[
\text{F at } y=y_0:\qquad
R_M(y_0)\bm c_m=0,\qquad
R_V(y_0)\bm c_m=0.
\]

因此，一切 L\'evy 型边界的\textbf{最终本征条件}都可写成
\[
\det K_m(k)=0,
\]
其中 \(K_m(k)\) 由 \(y=0,b\) 两边的四个边界行向量叠成一个 \(4\times4\) 矩阵。

下面列出最常见几类。

\paragraph{(a) SSSS}
\[
K_m^{\text{SS}}
=
\begin{pmatrix}
R_W(0)\\
R_M(0)\\
R_W(b)\\
R_M(b)
\end{pmatrix},
\qquad
\det K_m^{\text{SS}}=0.
\]
这一情形可进一步化简回 Navier 结果，最终频率为
\[
\omega_{mn}
=
\pi^2\sqrt{\frac{D}{\rho h}}
\left(
\frac{m^2}{a^2}+\frac{n^2}{b^2}
\right).
\]

\paragraph{(b) SSCC}
\[
K_m^{\text{SC}}
=
\begin{pmatrix}
R_W(0)\\
R_M(0)\\
R_W(b)\\
R_\theta(b)
\end{pmatrix},
\qquad
\det K_m^{\text{SC}}=0.
\]
这就是两条 \(x\)-边简支、两条 \(y\)-边固支时的最终频率方程。

\paragraph{(c) SSFF}
\[
K_m^{\text{SF}}
=
\begin{pmatrix}
R_W(0)\\
R_M(0)\\
R_M(b)\\
R_V(b)
\end{pmatrix},
\qquad
\det K_m^{\text{SF}}=0.
\]
这就是两条 \(x\)-边简支、两条 \(y\)-边自由时的最终频率方程。

\paragraph{(d) SSSC}
\[
K_m^{\text{S/C}}
=
\begin{pmatrix}
R_W(0)\\
R_M(0)\\
R_W(b)\\
R_\theta(b)
\end{pmatrix},
\qquad
\det K_m^{\text{S/C}}=0.
\]
这一式与 SSCC 同型；区别仅在于此处强调 \(y=0\) 为简支、\(y=b\) 为固支的非对称表述。

\paragraph{(e) SSSF}
\[
K_m^{\text{S/F}}
=
\begin{pmatrix}
R_W(0)\\
R_M(0)\\
R_M(b)\\
R_V(b)
\end{pmatrix},
\qquad
\det K_m^{\text{S/F}}=0.
\]

\paragraph{(f) SSCF}
\[
K_m^{\text{C/F}}
=
\begin{pmatrix}
R_W(0)\\
R_\theta(0)\\
R_M(b)\\
R_V(b)
\end{pmatrix},
\qquad
\det K_m^{\text{C/F}}=0.
\]

应当强调：对于 L\'evy 类问题，\textbf{本征频率的最终结果就是这些 \(4\times4\) 行列式为零所给出的超越方程}；一般不存在比这更短的统一闭式表达。

\subsubsection*{3.4 L\'evy 类静态问题}
静态时 \(k=0\)，故每个 \(m\) 阶满足
\[
\left(\frac{d^2}{dy^2}-\alpha_m^2\right)^2Y_m=\frac{q_m(y)}{D},
\qquad
q_m(y)=\frac{2}{a}\int_0^a q(x,y)\sin\frac{m\pi x}{a}\, dx.
\]
其齐次解退化为
\[
Y_m^{(h)}(y)
=
(A_m+B_my)\cosh(\alpha_my)+(C_m+D_my)\sinh(\alpha_my).
\]
因此静态 L\'evy 解的最终形式为
\[
w(x,y)=
\sum_{m=1}^\infty
\left[
Y_m^{(p)}(y)
+
(A_m+B_my)\cosh(\alpha_my)
+
(C_m+D_my)\sinh(\alpha_my)
\right]
\sin\frac{m\pi x}{a},
\]
其中 \(Y_m^{(p)}\) 是关于 \(q_m(y)\) 的任意特解，四个系数由 \(y=0,b\) 两边的边界条件确定。

\paragraph{均布载荷 \(q=q_0\)}
此时
\[
q_m(y)=\frac{4q_0}{m\pi},
\qquad m\ \text{为奇数};
\qquad
q_m(y)=0,\quad m\ \text{为偶数}.
\]
可取一个常数特解
\[
Y_m^{(p)}=\frac{4q_0}{D\,m\pi\,\alpha_m^4},
\qquad m\ \text{odd}.
\]
因此在所有 L\'evy 型静力问题中，最终解都归结为上述级数形式加上边界条件决定的四个常数。

\subsection*{4. 小结与组织方式}

以上内容可概括为以下四点：

\paragraph{1.}
在 Kirchhoff--Love 理论下，静态问题、自由振动、谐受迫振动本质上都是同一个线性算子
\[
D\nabla_t^4-\rho h\Omega^2
\]
的不同特例。

\paragraph{2.}
模态是理解一切问题的核心。只要求出本征对
\[
D\nabla_t^4\Phi_n=\rho h\omega_n^2\Phi_n,
\]
则静态解、自由振动解、谐响应都可以直接由模态展开写出。

\paragraph{3.}
圆板最自然的解析方法是极坐标分离变量：
\begin{itemize}
\item 静态问题对应双调和方程，最终结果由幂函数 / 对数函数给出；
\item 动态问题对应 Bessel 与修正 Bessel 函数；
\item C / S / F 三类边界条件分别给出不同的本征方程。
\end{itemize}

\paragraph{4.}
长方板最重要的解析方法有两类：
\begin{itemize}
\item \textbf{Navier 解}：四边简支，模态与频率完全显式；
\item \textbf{L\'evy 解}：一对对边简支，其余两边可为 C / S / F，最终结果由 \(4\times4\) 特征行列式给出。
\end{itemize}

\end{document}
